\chapter{Stato dell'arte e Background}

\label{capitolo2}

%\noindent Nella seconda sezione si riporta lo stato dell'arte del settore, un inquadramento dell'area di ricerca orientato a portare il lettore all'interno della problematica affrontata. Bisogna dimostrare di conoscere le cose fatte fino ad ora in questo campo e il perch\'e si sia reso necessario lo svolgimento di questo lavoro. Questa sezione deve essere grondante di citazioni bibliografiche~\cite{marco02bud}.

\section{Stato dell'arte}
All'interno delle reti sociali sta sempre pi\'u prendendo piede il problema della polarizzazione delle opinioni. Nel linguaggio comune il confronto tra individui ha sempre generato una forte controversia oppure un punto d'incontro. 
I social network hanno permesso all'utente di poter diffondere attraverso post, messaggi o espressioni audio video le proprie opinioni e/o pensieri all'interno di una comunità sociale. A tal proposito per favorire la diffusione delle diverse correnti di pensiero i social network stanno sempre più sviluppando algoritmi per permettere di identificare le comunità isolate, che condividono un unico punto di vista di un determinato argomento.
La polarizzazione è un algoritmo matematico che applicato all'interno delle rete sociali consente di capire quanto un utente, che accede per la prima volta all'interno di una rete sociale, venga influenzato dagli altri utenti e quanto una news o un giudizio si propaga all'interno di una rete sociale.
Prima di poter illustrare questo algoritmo, con le relative problematiche, verrà illustrata una definizione di rete sociale.

\subsection{Rete Sociale}
Una rete sociale consiste in un qualsiasi gruppo di individui connessi tra loro da diversi legami. Per gli esseri umani i legami vanno dalla conoscenza casuale, ai rapporti di lavoro, ai vincoli familiari. Le reti sociali sono spesso usate come base di studi interculturali in sociologia, in antropologia, in etologia.

L'analisi delle reti sociali, ovvero la mappatura e la misurazione delle reti sociali, può essere condotta con un formalismo matematico usando la teoria dei grafi. In generale, il corpus teorico ed i modelli usati per lo studio delle reti sociali sono compresi nella cosiddetta \textit{social network analysis}.

La ricerca condotta nell'ambito di diversi approcci disciplinari ha evidenziato come le reti sociali operino a pi\'u livelli e svolgano un ruolo cruciale nel determinare le modalità di risoluzione di problemi e i sistemi di gestione delle organizzazioni, nonché le possibilità dei singoli individui di raggiungere i propri obiettivi.

\subsection{Le reti}
La diffusione del web e del termine social network ha creato negli ultimi anni alcune ambiguità di significato. La rete sociale è infatti storicamente, in primo luogo, una rete fisica.

Rete sociale è, ad esempio, una comunità di lavoratori, che si incontra nei relativi circoli dopolavoristici e che costituisce una delle associazioni di promozione sociale. Esempi di reti sociali sono inoltre le comunità di sportivi, attivi o sostenitori di eventi, le comunità unite da problematiche strettamente lavorative e di tutela sindacale del diritto nel lavoro, le confraternite e in generale le comunità basate sulla pratica comune di una religione e il ritrovo in chiese, templi, moschee, sinagoghe e altri luoghi di culto.

La rete sociale ha una sua versione all'interno di \textbf{Internet} definita anche come \textit{Social media}, questa è una delle forme più evolute di comunicazione in rete. La rete delle relazioni sociali che ciascuno di noi tesse ogni giorno, in maniera più o meno casuale, nei vari ambiti della nostra vita, si può così "materializzare", organizzare in una "mappa" consultabile, e arricchire di nuovi contatti. I principali social network sono: \textit{Facebook, MySpace, Instagram, Twitter, Google+, LinkedIn, Ask.fm, Pinterest, Formspring, Bebo, Friendster, Hi5, Ning, Tagged, Meetup, Tumblr}.

\subsection{Twitter}
\textit{Twitter} è il social network di riferimento utilizzato per la realizzazione di questa tesi di laurea, ne diamo una breve descrizione, è una rete sociale, creata il 21 marzo 2006 dalla \textit{Obvious Corporation} di San Francisco, che fornisce agli utenti, attraverso l'omonima piattaforma, una pagina personale aggiornabile tramite messaggi di testo con lunghezza massima di 140 caratteri; nel 2017 l'azienda ha aumentato la lunghezza dei tweet a 280 caratteri per alcuni paesi. Gli aggiornamenti di stato possono essere effettuati tramite il sito stesso, via SMS, con programmi di messaggistica istantanea, posta elettronica, oppure tramite varie applicazioni basate sulle API di Twitter. 

Il nome "Twitter" deriva dal verbo inglese to tweet che significa "cinguettare". Tweet è anche il termine tecnico degli aggiornamenti del servizio. I Tweet che contengono esattamente 140 caratteri vengono chiamati \textit{Twoosh}. Gli aggiornamenti sono mostrati nella pagina di profilo dell'utente e comunicati agli utenti che si sono registrati per riceverli. È anche possibile limitare la visibilità dei propri messaggi oppure renderli visibili a chiunque. \begin{figure}[h!]
    \begin{center}
      \includegraphics[scale= 0.15]{twitter-logo.png}
	%\psfig{file=./pictures/logo}
	\caption{Logo di Twitter}
    \end{center}
  \end{figure}

Questo social network è molto utilizzato da personaggi famosi, sportivi, ma anche capi di stato, esponenti politici ed economici, perfino il Papa, consentendo a tutti gli utenti iscritti alla piattaforma di poter visualizzare le idee ed opinioni espresse da questi personaggi. Proprio per questo motivo \textit{Twitter} consente una diffusione su larga scala di idee ed opinioni da parte di moltissimi utenti, gli stessi giornalisti utilizzano questa piattaforma per diffondere \textit{scoop} o notizie di vario genere, proprio per la grande visibilità che fornisce questo social network.
Vengono generati moltissimi dati ed è proprio per questo motivo che ho deciso di utilizzarlo come strumento per la raccolta dati e cercare di calcolare la polarizzazione di alcuni topic selezionati, permettendomi di visualizzare come le informazioni espresse si propagano all'interno della rete e come queste influenzano il pensiero degli utenti.

\subsubsection{Tweet}
L'elemento alla base di questo social network è il \textit{Tweet}, che non è altro che un messaggio contenente le opinioni, le idee, i pensieri dell'utente. Tutte queste cose possono essere arricchite utilizzato elementi multimediali come video, immagini oppure link a pagine web. Il tweet una volta pubblicato potrà essere accessibile per tutti gli utenti delle rete che hanno una relazione di "amicizia" con l'utente in questione, potendo anche effettuare delle operazioni aggiuntive quale commentare la notizia, apprezzarla oppure \textit{retwettarla} cioè condividere queste informazioni a tutti i propri collegamenti.
I Tweet possono essere anche etichettati utilizzando due strumenti messi a disposizione da questo social network:
\begin{itemize}
\item \textit{Mentions}: sono dei riferimento ad utenti, cioè come assegnare un mittente ad un email, affermando a tutti i membri all'interno della rete, che quel tweet è destinato esclusivamente a quella persona. Può essere utile per condividere più informazioni con gli amici delle rete più stretti, oppure per formulare degli attacchi, riflessioni o considerazioni anche con utenti famosi. Per fare questa operazione è necessario inserire la \@ ed il nome dell'utente in questione per poter collegare al tweet in questione alla persona desiderata.
\item \textit{Hashtag}: parole o combinazioni di parole concatenate precedute dal simbolo cancelletto (\#).
Etichettando un messaggio con un hashtag si crea un collegamento ipertestuale a tutti i messaggi recenti che citano lo stesso hashtag. Molto  vengono utilizzati durante \textit{Show Tv}, \textit{comizi}, \textit{pubblicità}, \textit{eventi politici} per poter consentire alle persone che vogliano esprimere la loro opinioni sugli argomenti espressi durante questi eventi. 

\end{itemize}

\subsection{Polarizzazione}
Definiamo in maniera concettuale il significato, le funzionalità della polarizzazione. 
Citando il paper \textit{A Measure of Polarization on Social Media NetworksBased on Community Boundaries} di \textit{Pedro H. Calais Guerra, Wagner Meira Jr. Claire Cardie, Robert Kleinberg}(\cite{Guerra2015polarization} la polarizzazione viene definita come un processo sociale in cui un gruppo viene diviso: 
in due sottogruppi aventi un conflitto o una visione differente del problema in questione, e da alcuni individui che rimangono neutrali.
Comprendere la polarizzazione e quantificarla è una sfida a lungo termine per i ricercatori all'interno di diverse aree, proprio perché risulta essere un potente strumento per l'analisi delle opinioni.

Nel paper \textit{Measuring Political Polarization: Twitter shows the two sides of Venezuela} di \textit{A. J. Morales, J. Borondo, J. C. Losada e R. M. Benito} (\cite{morales2015polarization}), la polarizzazione assume una connotazione molto simile, vista in chiave politica, infatti all'interno di questo paper tale analisi viene effettuata all'interno di un contesto politico. 
La definizione che viene data della polarizzazione è: 
un fenomeno sociale che si verifica quando gli individui contrappongono le proprie credenze ed opinioni in una posizione conflittuale tra loro, mentre alcuni mantengono una posizione neutrale.
Per citare \textit{Jonh Turner}\footnote{Noto psicologo sociale britannico, con altri colleghi ha sviluppato la teoria dell'Auto-categorizzazione, che afferma tra l'altro che il sé non è un aspetto fondamentale della cognizione, ma è il risultato di processi cognitivi ed interazioni tra la persona e il contesto sociale.} :"Come le molecole polarizzate, i membri di un gruppo allineano il proprio pensiero nella direzione in cui erano originalmente diretti". 

Adottare questo strumento all'interno di un contesto politico fornisce una chiara visione della preferenza o meno verso i candidati politici, verificando il trend del candidato.
Dal punto di vista matematico possiamo definire la polarizzazione tra due gruppi come la distanza che intercorre tra due numeri. Definendo in questo modo quanto due idee siano diverse e contrastanti utilizzando la distanza tra due valori numerici di riferimento.
Le problematiche evidenziate dalla polarizzazione vengono trattate all'interno del paper : \textit{Reducing Controversy by Connecting Opposing Views} di \textit{Kiran Garimella ,Gianmarco De Francisci Morales, Aristides Gionis e Michael Mathioudakis} \cite{Garimella2015polarization}. Quando una popolazione si divide in due gruppi con visioni opposte quello che spesso si può verificare è la creazione degli \textit{Echo-chambers} cioè: una situazione dove persone che la pensano allo stesso modo rafforzano le proprie convinzioni a vicenda, senza esporsi verso idee contrastanti.
In conclusione la polarizzazione può aiutare ad individuare e a prevenire la formazioni degli \textit{Echo-chambers}.

\subsubsection{Echo-Chambers}
Brevemente illustriamo un gravissimo problema che abbiamo precedentemente riscontrato all'interno della polarizzazione.
Gli \textit{Echo-chambers} è una descrizione metaforica di una situazione in cui le informazioni, le idee, o credenze sono rinforzate dalla comunicazione e dalla ripetizione all'interno di un ambiente definito. Praticamente viene diffuso lo stesso messaggio, come un "\textit{eco}" che riproduce lo stesso suono più e più volte all'interno di un ambiente\cite{echo}. 
Nei social media il fenomeno degli echo-chambers è un problema piuttosto comune e ciò è dovuto dalla forte presenza di comunità di utenti che condividono sempre le stesse informazioni con un forte tasso di pubblicazione, quindi un nuovo utente che appoggia una di quelle idee entra all'interno di quel circolo vizioso. Ovviamente gli amministratori di questi social network sviluppano algoritmi che possano prevenire la formazione di queste comunità.
Le conseguenze di questo problema sono molto evidenti ovvero:
\begin{itemize}
\item \textit{Impossibilità della crescita culturale}
\item \textit{Impossibilità della crescita sociale}
\item \textit{Rendere la comunicazione poco democratica}
\end{itemize}
La polarizzazione come precedentemente accennato è uno strumento che può aiutare la prevenzione per la formazione di queste comunità.

\subsection{Sentiment Analysis}

La Sentiment Analysis è lo strumento utilizzato per poter catalogare le informazioni degli utenti attraverso un'analisi del testo. Non è possibile citare la sentiment analysis senza menzionare l' "opinion mining" ovvero una branchia del machine learning che studia il comportamento delle opinioni nella rete.
Il termine sentiment viene utilizzato in riferimento all'analisi automatica effettuata per valutare il testo ed esprimere un giudizio predittivo sul contenuto \cite{pang2015sentiment}.
Il compito base di questo processo è quello di classificare la polarità di un testo e quindi il suo contenuto e catalogarlo come \textit{positivo},\textit{negativo} o \textit{neutrale}.
Gli approcci più comunemente individuati rispetto alla sentiment analysis possono suddividersi in tre macro-categorie: 
\begin{itemize}
\item \textit{rilevamento delle keyword}: questo metodo consente di classificare il testo tramite categorie emotive facilmente riconoscibili, individuate in base alla presenza di parole emotive non ambigue, come felice, triste, e annoiato.
\item \textit{affinità lessicale}: questo metodo non rileva solo le keyword emotive, ma assegna anche a parole arbitrarie \textit{un'affinità probabile} a emozioni particolari. Rispetto alla prima metodologia vista, l'affinità lessicale consente di affinare la selezione e l'attribuzione della polarità.
\item \textit{metodi statistici}: questi metodi si basano su elementi di apprendimento automatico. Per misurare l'opinione nel contesto e trovare la caratteristica che è stata giudicata, sono usate le relazioni grammaticali delle parole utilizzate. Le relazioni di dipendenza grammaticale sono ottenute attraverso la scansione approfondita del testo. Il processo di apprendimento da parte della macchina (anche detto machine learning) non è immediato, devono infatti essere costruiti dei modelli che associano a diverse tipologie di commenti una polarità e se necessario ai fini dell'analisi anche un topic.
\end{itemize}

Quando parliamo della sentiment analysis non possiamo non esporre una delle grandi limitazioni offerte da questo strumento, cioè l'incapacità della macchina di poter comprendere concetti emotivi complessi come l'ironia.
Ad esempio un commento ad un ritardo nel volo:
\begin{quote}
"Il mio volo è in ritardo. Splendido!"
\end{quote}
verranno interpretati e classificati dalla macchina come post dalla polarità positiva mentre invece dovrebbero essere assegnate delle negatività. Quindi in un contesto lessicale pieno di ironia l'attendibilità sarà inferiore rispetto ad un documento con informazioni oggettive.
Per migliorare l'accuratezza della predizione è necessario l'intervento dell'uomo, aumentando il volume dei dati di riferimento per l'analisi, banalmente milioni di post, possono alleviare le preoccupazioni sull'affidabilità a livello granulare, ossia di un singolo post.
In conclusione la sentiment analysis è uno strumento che può essere molto utile all'interno dei social media in quanto può consentire di catalogare le opinioni degli utenti attraverso l'analisi dei contenuto pubblicati all'interno della rete. Questa potenzialità è ciò che mi ha spinto nell'adottare tale tecnica all'interno della mia tesi, garantendo una  classificazione dei diversi \textit{Tweet} collezionati.

\begin{comment}
La sentiment analysis può essere sviluppata seguendo 4 approcci differenti:
\begin{itemize}
\item \textit{Spotting di parole chiave}: il testo viene classificato da categorie influenti basate sulla presenza di parole influenti come contento, triste, impaurito, annoiato.
\item \textit{Affinità lessicale}: oltre a rilevare le parole influenzanti, assegna arbitrariamente alle parole una probabile affinità a emozioni particolari. 
\item \textit{Metodi statistici}:  questo approccio si basa sugli elementi trattati dal machine learning come analisi semantica latente, macchine a vettori di supporto, bag of words. Partendo da questi elementi per estrarre le opinioni in un contesto e ottenerne delle caratteristiche, vengono utilizzate delle relazioni grammaticali tra le parole. Le relazioni sono ottenute da un'analisi sintattica e profonda del testo.
\item \textit{Tecniche di livello concettuale}: utilizza gli elementi della rappresentazione della conoscenza come le \textit{ontologie} e le \textit{reti semantiche}, riuscendo a rilevare semantiche che sono espresse in maniera sottile all'interno del testo.
\end{itemize}
In conclusione la sentiment analysis è uno strumento che può essere molto utile all'interno dei social media in quanto può consentire di catalogare le opinioni degli utenti attraverso l'analisi dei contenuto pubblicati all'interno della rete. Questa potenzialità è ciò che mi ha spinto nell'adottare tale tecnica all'interno della mia tesi.
\end{comment}

\subsection{Forecasting}
Il Forecasting è un processo per fare predizioni del futuro basandosi sui dati passati e presenti mediante l'analisi delle tendenze. Questo viene utilizzato molto all'interno delle aziende per poter analizzare i guadagni ed i futuri ricavi. \cite{Forecasting}
Spesso viene associato ai processi di budgeting e pianificazione, si serve di dati passati e presenti, analisi dei trend e informazioni esecutive per prevedere la situazione futura di ogni  indicatore.
Il forecasting prevede tre diverse tecniche:
\begin{itemize}
\item \textit{Qualitativa}
\item \textit{Analisi e proiezione delle serie temporali}
\item \textit{Modelli causali}
\end{itemize}
La previsione finanziaria non è una scienza esatta e l'incertezza ne è un tipico aspetto. Per limitare gli errori e migliorare l'esattezza delle previsioni è necessario utilizzare precisi dati storici e in tempo reale. Inoltre, la possibilità di combinare previsioni rolling, previsioni a lungo raggio, simulazione di scenari possibili e stress test supporta i risultati delle previsioni e rende le previsioni più agili e reattive ai cambiamenti economici e di business.
Ai fini della mia tesi ho deciso di adottare uno studio basato sulle analisi e proiezioni delle serie temporali, perché permette di estrarre statistiche significative ed altre caratteristiche dei dati. Predice i valori futuri attraverso l'osservazione dei dati precedentemente raccolti.




